\scanpage{7}
% 1–3
\uline{Cor} $A_1\subset S_1(k)$ stable par primitivation canonique (i.e. celle qui s'annule à l'origine), contient $\log(1+t_1)$, $\exp t_1$, $\cos t_1$, $\sin t_1$, $\mathrm{tg}\,t_1$ ....

% 4–8
Soit $X$ schéma formel lisse de dim $n=p+q$. On sait \deleted{que} ce que c'est qu'un \uline{feuilletage} \inserted{de dim $p$} de $X$ : variété formelle quotient $B$ de $X$, avec $X-B$ submersif $\}$, $\dim B=q=n-p$. A un tel feuilletage est associé un sous-module $D\subset\Theta(X)$

% 9–12
[champs de vecteurs le long de la fibre] qui est un ``système de Pfaff de dim $p$ complètement intégrable'', et vice-versa. Si $X$ a une $A$-structure, on dit qu'on a un $A$-feuilletage si $\exists$ \inserted{sur} $B$ une $A$-structure

% 13–17
(d'ailleurs unique \unsure{$\}$|)|,} grâce à f)) telle que $X\to B$ soit un $A$-morphisme. Cela \uline{implique} que le sous-module $D\subset\Theta(X)$ (ou sous-fibré de $T(X)$) soit un $A$-sous-fibré -- i.e. qu'on ait un $A$-syst. de Pfaff. compl. intégrable. On veut qu'on ait une réciproque, savoir

% 18–20
h) Si $D_A$ est un $A$-syst. de Pfaff compl. intégrable sur le $A$-schéma formel $X$, alors il donne lieu à un $A$-feuilletage \deleted{\unsure{complèt$^{\mathrm{t}}$ intégr.|complét. intégr.|réalisé}}. [En d'autres termes,

% 18, left
% slanted
\margin{\uline{NB} h) implique g), qui équivaut à h) pour $p=1$}

% 21–27
si $i:Y-X$ est un $A$-morphisme donné tel que le composé
\[
 T(Y)_0\mathrel{\overset{T(i)_0}{\text{---}}}T(X)_0-T(X)_0/D(0)
\]
soit un isom, alors l'unique morphisme
\[
 \pi:X-Y
\]
tel que $\pi\circ i=\mathrm{id}_Y$ et que $\varphi$ soit compatible avec $D$, est un $A$-morphisme ....

% 28, foot
\margin{$\Longleftarrow$} h') la fibre de $0$ dans un tel feuilletage est une sous-$A$-variété formelle ....]
% archive sheet number: 6
